\pdfoutput=1
\documentclass[twocolumn]{aastex631}

\usepackage{amsmath}
\usepackage{amssymb}
\usepackage{graphicx}
\usepackage{booktabs}

\input{numbers.tex}

\newcommand{\Ha}{\ensuremath{\mathrm{H}\alpha}}
\newcommand{\Hb}{\ensuremath{\mathrm{H}\beta}}
\newcommand{\OIII}{\ensuremath{[\mathrm{O\,\textsc{iii}}]}}
\newcommand{\NII}{\ensuremath{[\mathrm{N\,\textsc{ii}}]}}
\newcommand{\SII}{\ensuremath{[\mathrm{S\,\textsc{ii}}]}}
\newcommand{\OI}{\ensuremath{[\mathrm{O\,\textsc{i}}]}}
\newcommand{\HeI}{\ensuremath{\mathrm{He\,\textsc{i}}}}
\newcommand{\kms}{\ensuremath{\mathrm{km\,s^{-1}}}}
\newcommand{\LHa}{\ensuremath{L_{\mathrm{H}\alpha}}}
\newcommand{\FWHM}{\ensuremath{\mathrm{FWHM}}}
\newcommand{\FWHMsg}{\ensuremath{\mathrm{FWHM}_{\mathrm{sg}}}}
\newcommand{\fvir}{\ensuremath{f_{\mathrm{vir}}}}
\newcommand{\MBH}{\ensuremath{M_{\mathrm{BH}}}}
\newcommand{\sigV}{\ensuremath{\sigma(V)}}

\shorttitle{Velocity-Window Dependence of the LRD \Ha\ Relation}
\shortauthors{Wang}
\hypersetup{pdftitle={Velocity-Window Dependence of the Broad H-alpha Width-Luminosity Relation in Little Red Dots},pdfauthor={Haoyu Wang}}

\begin{document}

\title{Velocity-Window Dependence of the Broad \texorpdfstring{\Ha}{H-alpha} Width--Luminosity Relation in Little Red Dots}

\correspondingauthor{Haoyu Wang}
\email{pkuwanghaoyu@pku.edu.cn}

\author[0009-0005-9117-0842]{Haoyu Wang}
\affiliation{Independent Researcher}

\begin{abstract}
Broad Balmer emission lines probe black hole masses and nuclear kinematics in high-redshift active galactic nuclei (AGNs). In JWST ``little red dots'' (LRDs), integrated broad-\Ha\ widths span a remarkably narrow range at fixed line luminosity, raising the question of whether this tightness persists once outer-wing emission is included. Holding each object and its broad-\Ha\ luminosity fixed, we measure the finite-window line dispersion \sigV, the flux-weighted second-moment width within $\pm V$, from $V = 1500$ to $4500~\kms$ in JWST/NIRSpec spectra and Sloan Digital Sky Survey (SDSS) Type-1 AGNs. Across $\PTJWSTPresetN$ high signal-to-noise LRDs, the observed scatter in $\log_{10}\sigV$ at fixed \LHa\ rises from $\PTJWSTPresetSTwoFiveThreeDec$ to $\PTJWSTPresetSFourFiveThreeDec$~dex between $V = 2500$ and $4500~\kms$, a paired increase of $\PTJWSTPresetDeltaThreeDec$~dex with a 95\% bootstrap interval of $\PTJWSTPresetDeltaLoThreeDec$--$\PTJWSTPresetDeltaHiThreeDec$~dex, compared with $\PTSDSSPresetDeltaThreeDec$~dex across $\PTSDSSPresetNFourFive$ SDSS AGNs. Injecting synthetic profiles into the empirical noise of the same $\PTJWSTPresetN$ spectra shows that single-Gaussian and core-plus-wing models both accommodate this rise, leaving intrinsic wing diversity, truncation, and wing noise unseparated. In the clearest dual-line spectrum, RUBIES-EGS-55604, \Hb\ strengthens relative to \Ha\ from line center to wing under two spectral decompositions of the same grating spectrum. Alongside integrated widths in JWST, SDSS, and $\EuclidMainN$ Euclid DR1 AGNs, these results show that the tight \Ha\ scaling depends on the velocities included in \sigV, requiring mass and kinematic models to match both transition and velocity interval.
\end{abstract}

\keywords{Active galactic nuclei (16) --- Quasars (1319) --- Emission line galaxies (459) --- High-redshift galaxies (734) --- Supermassive black holes (1663)}

\section{Introduction}
\label{sec:intro}

JWST spectroscopy has revealed an abundant population of compact rest-optical sources at $z \gtrsim 3$, termed ``little red dots'' (LRDs), with prominent broad Balmer emission lines \citep{Matthee2024,Greene2024} and characteristic V-shaped continua inflecting near the Balmer limit \citep{Setton2025Vshape}. In several bright examples with simultaneous \Ha\ and \Hb\ grating coverage, \citet[][Section~V.3.1 and Figure~9]{Chen2026} further showed that the $F(\Ha)/F(\Hb)$ flux ratio decreases substantially from line center toward the wings, indicating that emission or radiative-transfer conditions vary across an individual broad-line profile.

Whether LRD broad lines trace virialized motion around massive black holes or radiative transfer in dense gas envelopes remains actively debated \citep[e.g.,][]{Rusakov2026,Madau2026b,InayoshiHo}. In classical active galactic nuclei (AGNs), single-epoch black hole mass estimates follow $\MBH = \fvir\, R\, \Delta V^2 / G$ \citep{Peterson2004,GreeneHo2005,Bentz2013}, where the virial factor \fvir\ depends on whether $\Delta V$ is characterized by the full width at half maximum (\FWHM) or the flux-weighted line dispersion \citep{Collin2006}. Across AGN populations, width scatter at fixed luminosity constrains broad-line geometry and the dynamical information carried by single-epoch widths \citep{Fine2008,Croom2011,Assef2012}, while optical-depth effects and scattering can reshape emergent line wings relative to bulk gas velocities \citep{OsterbrockFerland2006,Laor2006}.

Whereas local quasars typically show $\gtrsim 0.12$~dex of width scatter at fixed luminosity, \citet{Wang2026} found that JWST LRDs follow a tight broad-\Ha\ width--luminosity relation with a measurement-error-corrected intrinsic scatter of $0.054$~dex at fixed \LHa, less than half that of refitted SDSS quasars and independently seen in \citet{BaccusXu2026}. \citet{Wang2026} also found similarly small scatter for single-Gaussian \FWHM\ (\FWHMsg), interpercentile width $W_{50}$, and Lorentzian or exponential fits, alongside larger scatter for two-Gaussian composite \FWHM. Because full-profile summary widths weight the bright core heavily whereas scattering, outflows, or multi-component kinematics act most strongly in the fainter wings, changing the summary statistic does not isolate individual velocity ranges. Here we hold the objects, width definition, and broad-\Ha\ luminosities fixed, varying only the velocity range $[-V, +V]$ in the finite-window line dispersion \sigV\ to test whether LRD widths remain tightly clustered once outer wings are included.

Section~\ref{sec:res_moments} and Figure~\ref{fig:moments} measure how the observed \Ha\ scatter changes across $V = 1500$--$4500~\kms$ in JWST LRDs and SDSS Type-1 AGNs, calibrated by empirical-noise profile injections on the same targets. Section~\ref{sec:res_balmer} and Figure~\ref{fig:balmer} examine velocity-binned $F(\Hb)/F(\Ha)$ ratios under two decompositions of the same combined grating spectra in three high-weight dual-line LRDs and relate them to a sample-wide template fit. Section~\ref{sec:res_context} and Figure~\ref{fig:three_survey} place integrated LRD \Ha\ widths alongside SDSS and Euclid DR1 AGNs, and Section~\ref{sec:discussion} discusses implications for broad-line modeling.

\section{Observations and Samples}
\label{sec:samples}

Table~\ref{tab:samples} summarizes all spectroscopic samples and subsets analyzed here. Because the three surveys span disjoint regimes of redshift, spectral resolution, and continuum luminosity, each serves a distinct empirical role rather than forming a fixed-luminosity evolutionary sequence.

\begin{deluxetable*}{lcccp{3.75cm}}
\tabletypesize{\footnotesize}
\tabcolsep=3.8pt
\tablewidth{0pt}
\tablecaption{Samples and Measurements\label{tab:samples}}
\tablehead{
\colhead{Survey / Subset} &
\colhead{$N$} &
\colhead{Redshift} &
\colhead{Measured Quantities} &
\colhead{Selection and Analytical Role}
}
\startdata
\cutinhead{JWST/NIRSpec Little Red Dots (Archival Grating Parent: $N = \JWSTGratingParentN$)}
Continuum-clean baseline &
  $\JWSTCleanN$\tablenotemark{a} &
  $\JWSTCleanZMin$--$\JWSTCleanZMax$ &
  \Ha\ \FWHMsg &
  Continuum-verified integrated \Ha\ width--luminosity baseline \\
Primary window sample &
  $\PTJWSTPresetN$\tablenotemark{b} &
  $\JWSTPresetZMin$--$\JWSTPresetZMax$ &
  \Ha\ \sigV, \FWHMsg &
  Broad \Ha\ $\mathrm{S/N} \ge 15$ with both default continuum sidebands \\
Supplementary window sample &
  $\PTJWSTFullN$\tablenotemark{b} &
  $\JWSTHighSNRZMin$--$\JWSTHighSNRZMax$ &
  \Ha\ \sigV, \FWHMsg &
  Adds $\JWSTEdgeFallbackN$ detector-edge sources via near-line red continuum sideband \\
Same-grating Balmer sample &
  $\POneTmplAllThirtyThreeN$ (top $\POneTmplTopThreeN$)\tablenotemark{c} &
  $\JWSTTwoLineZMin$--$\JWSTTwoLineZMax$ &
  \Hb/\Ha\ $R(v)$, template &
  Simultaneous \Ha\ and \Hb\ on same combined grating spectrum \\
\cutinhead{Comparison and Population-Context Surveys}
SDSS DR7 Type-1 AGNs &
  $\SDSSCleanN$ ($\PTSDSSPresetNFourFive$--$\PTSDSSPresetNOneFive$)\tablenotemark{d} &
  $\SDSSCleanZMin$--$\SDSSCleanZMax$ &
  \Ha\ \sigV, \FWHMsg &
  Same-window empirical reference in local broad-line AGNs \\
Euclid DR1 primary sample &
  $\EuclidMainN$\tablenotemark{e} &
  $\EuclidMainZMin$--$\EuclidMainZMax$ &
  \Ha\ \FWHMsg, $\sigma_w$ &
  Independently selected intermediate-$z$ integrated-width reference \\
\enddata
\tablenotetext{a}{Subset of the $\JWSTGratingParentN$ archival NIRSpec grating LRDs passing local continuum verification around \Ha.}
\tablenotetext{b}{All $\PTJWSTFullN$ continuum-clean sources with broad-\Ha\ $\mathrm{S/N} \ge 15$, comprising the $\PTJWSTPresetN$ primary default-sideband sources plus $\JWSTEdgeFallbackN$ detector-edge objects.}
\tablenotetext{c}{Selected directly from the $\JWSTGratingParentN$ grating sources, comprising $\JWSTTwoLineInCleanN$ of the $\JWSTCleanN$ continuum-clean objects and $\JWSTTwoLineOutCleanN$ objects outside that subset; velocity-binned ratios are analyzed in the $\POneTmplTopThreeN$ highest-weight objects.}
\tablenotetext{d}{From $\SDSSSpectra$ spectra cross-matched to \citet{Shen2011}, $\SDSSCleanN$ pass single-Gaussian quality control and $\PTSDSSPresetNFourFive$--$\PTSDSSPresetNOneFive$ yield valid moments across $V = 1500$--$4500~\kms$ under default sidebands, rising to $\PTSDSSFullNFourFive$--$\PTSDSSFullNOneFive$ with near-line sidebands; see Table~\ref{tab:moments}.}
\tablenotetext{e}{Selected from $\EuclidTargets$ Euclid DR1 candidates at $19.0 < H_{\mathrm{aper}} < \EuclidHLimit$ with usable MER photometry prior to inspecting broad-line widths.}
\end{deluxetable*}

\emph{JWST/NIRSpec LRDs.}---We start from an archival parent compilation of $\JWSTGratingParentN$ deduplicated LRDs observed with JWST/NIRSpec medium- or high-resolution gratings at $R \sim 1000$--$2700$ \citep{Jakobsen2022,deGraaffLSF2024}. Reduced 1D spectra were retrieved from the DAWN JWST Archive,\footnote{\url{https://dawn-cph.github.io/dja/}} extracted with \texttt{msaexp} \citep{msaexp} and including public observations from the RUBIES survey \citep{DJAdeGraaff2024}. Our target list combines published spectroscopic LRD compilations, primarily \citet{BaccusXu2026} and Table~1 of \citet{Yanagisawa2026} alongside \citet{Kocevski2025}, \citet{Taylor2025}, and \citet{Hviding2025}, retaining any source classified as an LRD in at least one contributing catalog. Cross-matching against the full JWST spectroscopic sample of \citet{Wang2026} yields $\JWSTGratingParentWangOverlapN$ shared sources among the $\JWSTGratingParentN$ parent targets, corresponding to $\JWSTCleanWangOverlapN$ of the $\JWSTCleanN$ continuum-clean sources, $\JWSTHighSNRWangOverlapN$ of the $\JWSTHighSNRN$ high-S/N sources, $\JWSTPresetWangOverlapN$ of the $\PTJWSTPresetN$ primary window sources, and $\JWSTTwoLineWangOverlapN$ of the $\POneTmplAllThirtyThreeN$ dual-line sources; these overlap counts include CEERS-1244, which served as a comparison AGN in \citet{Wang2026} but is classified as an LRD in Table~1 of \citet{Yue2024}. The $\JWSTGratingParentN$ grating spectra enter two analyses:
\begin{enumerate}
\item \emph{\Ha\ velocity-window sample.}---Of the $\JWSTGratingParentN$ grating sources, $\JWSTCleanN$ at $z = \JWSTCleanZMin$--$\JWSTCleanZMax$ (median $\JWSTCleanZMed$, including $\JWSTCleanBXProvN$ in \citealt{BaccusXu2026} and $\JWSTCleanAtlasProvN$ in \citealt{Yanagisawa2026}) pass local continuum verification for the integrated-width baseline. Single-Gaussian fits define the uniform comparison width \FWHMsg\ with median $\JWSTCleanMedFWHM~\kms$, whereas broad-\Ha\ luminosities, spanning $\log_{10}(\LHa/\mathrm{erg\,s^{-1}}) = \JWSTCleanLHaMin$--$\JWSTCleanLHaMax$ with median $\JWSTCleanLHaMed$, and broad-line S/N use the total broad flux selected by the Bayesian information criterion (BIC); a symmetric exponential or summed two-Gaussian profile is adopted only when it yields the lowest BIC by $\Delta\mathrm{BIC}>6$ relative to the single Gaussian, and each object's luminosity is held fixed across all velocity windows. Within these $\JWSTCleanN$ objects, $\JWSTHighSNRN$ satisfy $\mathrm{S/N} \ge 15$ across $z = \JWSTHighSNRZMin$--$\JWSTHighSNRZMax$ with median $\FWHMsg = \JWSTHighSNRMedFWHM~\kms$. Both default continuum sidebands are covered in $\PTJWSTPresetN$ of these $\JWSTHighSNRN$ spectra, forming our \emph{primary} window sample, while adding $\JWSTEdgeFallbackN$ detector-edge targets, RUBIES-EGS-50052 and JADES-GS-38562, via a near-line red continuum sideband defines the $\PTJWSTFullN$-source \emph{supplementary} sample detailed in Appendix~\ref{app:edge}.
\item \emph{Dual-line Balmer sample.}---Separately, $\POneTmplAllThirtyThreeN$ of the $\JWSTGratingParentN$ grating sources at $z = \JWSTTwoLineZMin$--$\JWSTTwoLineZMax$ cover \Ha\ and \Hb\ simultaneously on the same combined grating spectrum, comprising $\JWSTTwoLineInCleanN$ of the $\JWSTCleanN$ continuum-clean sources and $\JWSTTwoLineOutCleanN$ additional targets, RUBIES-UDS-149298, Skyfire-1090, and RUBIES-EGS-42803. Velocity-binned Balmer ratios are examined in the $\POneTmplTopThreeN$ objects carrying the highest inverse-variance weight in the $\POneTmplAllThirtyThreeN$-source two-line template fit: RUBIES-EGS-55604, RUBIES-EGS-42046, and JADES-GN-68797.
\end{enumerate}

\emph{SDSS DR7 Type-1 AGNs.}---As a low-redshift empirical reference measured with the identical finite-window pixel integration, we use $\SDSSCleanN$ quality-controlled broad-line AGNs at $z = \SDSSCleanZMin$--$\SDSSCleanZMax$ (median $\SDSSCleanZMed$, median $\FWHMsg = \SDSSCleanMedFWHM~\kms$) refitted from $\SDSSSpectra$ SDSS DR7 spectra cross-matched to \citet{Shen2011}. We convert the cataloged broad-\Ha\ luminosities of \citet{Shen2011} to our adopted cosmology, spanning $\log_{10}(\LHa/\mathrm{erg\,s^{-1}}) = \SDSSCleanLHaMin$--$\SDSSCleanLHaMax$ with median $\SDSSCleanLHaMed$. Default continuum sidebands yield $\PTSDSSPresetNFourFive$--$\PTSDSSPresetNOneFive$ valid spectra across $V = 1500$--$4500~\kms$, increasing to $\PTSDSSFullNFourFive$--$\PTSDSSFullNOneFive$ when near-line sidebands are allowed for edge spectra (Table~\ref{tab:moments}).

\emph{Euclid DR1 Primary Sample and Cosmology.}---To establish an independently selected intermediate-redshift reference, we use $\EuclidMainN$ broad-\Ha\ AGNs at $z = \EuclidMainZMin$--$\EuclidMainZMax$ (median $\EuclidMainZMed$, median $\FWHMsg = \EuclidMainMedFWHM~\kms$, and $\log_{10}(\LHa/\mathrm{erg\,s^{-1}}) = \EuclidMainLHaMin$--$\EuclidMainLHaMax$) measured from Euclid DR1 NISP slitless spectroscopy \citep{EuclidAussel2026,EuclidLeBrun2026,EuclidFumana2026}. We selected the parent list of $\EuclidTargets$ candidates at $19.0 < H_{\mathrm{aper}} < \EuclidHLimit$ with usable MER photometry prior to inspecting the broad-line widths, distinct from photometric LRD selections in Euclid Quick Data Release~1 \citep{EuclidBisigello2025}. Because NISP slitless resolution and per-pixel S/N do not support native-pixel wing moments out to $\pm 4500~\kms$, Euclid enters solely through integrated-width statistics. Across the shared luminosity band $\log_{10}(\LHa/\mathrm{erg\,s^{-1}}) = 43.0$--$43.7$, SDSS and Euclid contribute $\SDSSBridgeN$ and $\EuclidBridgeN$ objects whereas JWST contributes only $\JWSTBridgeN$, and the brightest JWST overlap continuum at $\log_{10}(\lambda L_\lambda/\mathrm{erg\,s^{-1}}) = \JWSTBridgeCMax$ lies $\BridgeJWSTEuclidContinuumGap$~dex below the faintest Euclid overlap source at $\EuclidBridgeCMin$. Throughout this paper, luminosity distances use the \texttt{Planck18} flat $\Lambda$CDM cosmology in Astropy with $H_0 = 67.66~\kms\,\mathrm{Mpc}^{-1}$ and $\Omega_{\mathrm{m},0} = 0.30966$ \citep{Planck2020}.

\section{Measurements and Analysis}
\label{sec:methods}

\subsection{Finite-Window Line Dispersion}
\label{sec:methods_moments}

For each JWST and SDSS spectrum in the \Ha\ moment scan, rest-frame velocity $v$ is referenced to the fitted broad-\Ha\ center. A linear inverse-variance continuum is fitted across default two-sided sidebands at $[-11000, -6500]$ and $[+8700, +11000]~\kms$, masking pixels within $450~\kms$ of rest-frame $\OI\,\lambda 6365$ and avoiding \SII. Within $|v| \le 4500~\kms$, a two-Gaussian broad model plus narrow \Ha\ and $\NII\,\lambda\lambda 6548,6583$ are fitted simultaneously solely to subtract the narrow lines. We retain the signed narrow-subtracted residual flux $f(v)$ on the native pixel grid, preserving Balmer absorption troughs and each instrument's native line-spread function (LSF), and center the moment on the flux centroid $\bar{v}_V$ within each symmetric window $[-V, +V]$:
\begin{equation}
\sigma^2(V) = \frac{\int_{-V}^{V} (v - \bar{v}_V)^2 f(v)\,dv}{\int_{-V}^{V} f(v)\,dv},
\quad
\bar{v}_V = \frac{\int_{-V}^{V} v\, f(v)\,dv}{\int_{-V}^{V} f(v)\,dv}.
\label{eq:sigma_v}
\end{equation}
We refer to $\sigV \equiv \sqrt{\sigma^2(V)}$ throughout as the \emph{finite-window line dispersion}. Exact positive-overlap pixel quadrature formulas are given in Appendix~\ref{app:exact_int}; to avoid bridging wavelength gaps, we cap each pixel half-width at $0.55$ times the median pixel spacing and require $\ge 98\%$ valid velocity coverage, positive integrated flux, and $\sigma^2(V) > 0$.

Two properties govern \sigV\ across $V \in \{1500, 2500, 3500, 4500\}~\kms$. First, for a nearly flat flux density inside a narrow window, Equation~(\ref{eq:sigma_v}) approaches $\sigV = V/\sqrt{3}$ as derived in Appendix~\ref{app:exact_int}, compressing differences between broader and narrower profiles; in addition, a fixed velocity cut $V$ reaches further into the wings of narrower LRD profiles, which have median $\FWHMsg = \JWSTCleanMedFWHM~\kms$, than of broader SDSS profiles with median $\SDSSCleanMedFWHM~\kms$. Second, the $(v - \bar{v}_V)^2$ factor amplifies weak-wing noise and continuum-placement errors at large $V$, a sensitivity familiar from repeated single-epoch AGN spectroscopy \citep{Denney2009}. In empirical four-pixel block-noise recovery tests defined prior to measuring the observed moments (Appendix~\ref{app:injections}), high-S/N JWST spectra satisfy all qualification thresholds through $V = 4500~\kms$, achieving 68th-percentile recovery errors of $\PTCalJWSTGaussFourFiveErrSixEight$ and $\PTCalJWSTWingFourFiveErrSixEight$~dex for Gaussian and core-plus-wing inputs, respectively, before losing joint qualification at $V = 5500$ and $6000~\kms$.

\subsection{Velocity-Binned Balmer Flux Ratios}
\label{sec:methods_balmer}

We measure resolved Balmer profiles in RUBIES-EGS-55604 at $z_{\OIII} = \POneSrcFiveFiveZOIII$ (G395M, $\POneSrcFiveFiveNCombine$ combined exposures), RUBIES-EGS-42046 at $z_{\OIII} = \POneSrcFourTwoZOIII$ (G395M, $\POneSrcFourTwoNCombine$ combined exposures), and JADES-GN-68797 at $z_{\OIII} = \POneSrcSixEightZOIII$ (G395H, $\POneSrcSixEightNCombine$ combined exposures). Velocities are referenced to local systemic \OIII, rather than the broad-\Ha\ centroid used in Section~\ref{sec:methods_moments} and in the $\POneTmplAllThirtyThreeN$-source two-line template, and native-pixel overlaps are integrated across five symmetric $|v|$ intervals: $0$--$300$, $300$--$1000$, $1000$--$1500$, $1500$--$2500$, and $2500$--$3500~\kms$. On each combined grating spectrum, we evaluate $R \equiv F(\Hb)/F(\Ha)$ both after narrow-line subtraction and LSF matching, denoted decomposition~A, and after subtracting only a local linear continuum around each line, denoted decomposition~B. Both decompositions retain Balmer absorption without infill; equivalent Balmer decrements $D = 1/R$ are quoted only where integrated $\Hb \ge 2\sigma$.

\subsection{Residual Scatter and Resampling}
\label{sec:methods_regression}

At each window half-width $V$, we report $s(V)$ in dex as the sample standard deviation of the residuals about a separate ordinary least-squares (OLS) linear regression of $\log_{10}\sigV$, or $\log_{10}\FWHMsg$ for integrated widths, on $\log_{10}\LHa$. Single-window uncertainties resample valid sources with replacement and refit the relation on each draw. For the window-to-window change $\Delta s_{2500\to 4500} \equiv s(4500) - s(2500)$, we resample objects measured in both windows as pairs and refit both relations on every bootstrap draw, whereas cross-survey scatter differences resample each survey independently. Throughout, $s(V)$ and $\Delta s_{2500\to 4500}$ measure \emph{observed} residual scatter rather than measurement-error-corrected intrinsic scatter, and \LHa-conditioned scatter is kept distinct from Euclid's continuum-conditioned, selection-corrected intrinsic scatter $\sigma_w$.

\section{Results}
\label{sec:results}

\subsection{Velocity-Window Response of the Observed \texorpdfstring{\Ha}{H-alpha} Scatter}
\label{sec:res_moments}

\begin{figure*}[t]
\centering
\includegraphics[width=\textwidth]{figures/fig1_cumulative_dispersion.pdf}
\caption{Observed scatter in broad-\Ha\ finite-window line dispersion \sigV\ at fixed \LHa\ under unified default continuum sidebands. \textbf{(a)} Per-window OLS residual scatter $s(V)$ across the validated range $V = 1500$--$4500~\kms$ for the $\PTJWSTPresetN$ primary JWST LRDs (crimson circles and solid curve) and $\PTSDSSPresetNFourFive$--$\PTSDSSPresetNOneFive$ SDSS Type-1 AGNs (steel-blue squares and dashed curve); shaded bands denote per-window 95\% source-bootstrap confidence intervals, with extended-window diagnostics in Appendix~\ref{app:calibration}. \textbf{(b)} Paired scatter increment $\Delta s_{2500\to 4500}$ from $V = 2500$ to $4500~\kms$. Upper rows show observed paired increments with 95\% source-bootstrap confidence intervals (filled symbols with capped solid bars) for the $\PTJWSTPresetN$ JWST sources and the $\PTSDSSPresetNFourFive$ SDSS sources valid in both windows. Lower rows show matching 95\% conditional prediction intervals (capless shaded dashed bands with open median markers) from synthetic profile injections into the empirical noise of the same $\PTJWSTPresetN$ JWST spectra for single-Gaussian and core-plus-wing models; the vertical dotted crimson line marks the observed $\PTJWSTPresetN$-source JWST increment.}
\label{fig:moments}
\end{figure*}

\begin{deluxetable*}{lccc}
\tabletypesize{\footnotesize}
\tablewidth{0pt}
\tablecaption{Observed Residual Scatters $s(V)$ and Paired Increments Across Validated Velocity Windows\label{tab:moments}}
\tablehead{
\colhead{Window Half-Width / Quantity} &
\colhead{JWST (dex)} &
\colhead{SDSS (dex)} &
\colhead{SDSS $N$}
}
\startdata
\cutinhead{Primary Comparison: Default Continuum Sidebands (JWST $N = \PTJWSTPresetN$)}
$V = 1500~\kms$ &
  $\PTJWSTPresetSOneFive$ $[\PTJWSTPresetSLoOneFive, \PTJWSTPresetSHiOneFive]$ &
  $\PTSDSSPresetSOneFive$ $[\PTSDSSPresetSLoOneFive, \PTSDSSPresetSHiOneFive]$ &
  $\PTSDSSPresetNOneFive$ \\
$V = 2500~\kms$ &
  $\PTJWSTPresetSTwoFive$ $[\PTJWSTPresetSLoTwoFive, \PTJWSTPresetSHiTwoFive]$ &
  $\PTSDSSPresetSTwoFive$ $[\PTSDSSPresetSLoTwoFive, \PTSDSSPresetSHiTwoFive]$ &
  $\PTSDSSPresetNTwoFive$ \\
$V = 3500~\kms$ &
  $\PTJWSTPresetSThreeFive$ $[\PTJWSTPresetSLoThreeFive, \PTJWSTPresetSHiThreeFive]$ &
  $\PTSDSSPresetSThreeFive$ $[\PTSDSSPresetSLoThreeFive, \PTSDSSPresetSHiThreeFive]$ &
  $\PTSDSSPresetNThreeFive$ \\
$V = 4500~\kms$ &
  $\PTJWSTPresetSFourFive$ $[\PTJWSTPresetSLoFourFive, \PTJWSTPresetSHiFourFive]$ &
  $\PTSDSSPresetSFourFive$ $[\PTSDSSPresetSLoFourFive, \PTSDSSPresetSHiFourFive]$ &
  $\PTSDSSPresetNFourFive$ \\
Paired $\Delta s_{2500\to 4500}$ &
  $\PTJWSTPresetDeltaSigned$ $[\PTJWSTPresetDeltaLo, \PTJWSTPresetDeltaHi]$ &
  $\PTSDSSPresetDeltaSigned$ $[\PTSDSSPresetDeltaLo, \PTSDSSPresetDeltaHi]$ &
  $\PTSDSSPresetNFourFive$ \\
\cutinhead{Supplementary Check: Near-Line Sidebands Allowed for Edge Spectra (JWST $N = \PTJWSTFullN$)}
$V = 1500~\kms$ &
  $\PTJWSTFullSOneFive$ $[\PTJWSTFullSLoOneFive, \PTJWSTFullSHiOneFive]$ &
  $\PTSDSSFullSOneFive$ $[\PTSDSSFullSLoOneFive, \PTSDSSFullSHiOneFive]$ &
  $\PTSDSSFullNOneFive$ \\
$V = 2500~\kms$ &
  $\PTJWSTFullSTwoFive$ $[\PTJWSTFullSLoTwoFive, \PTJWSTFullSHiTwoFive]$ &
  $\PTSDSSFullSTwoFive$ $[\PTSDSSFullSLoTwoFive, \PTSDSSFullSHiTwoFive]$ &
  $\PTSDSSFullNTwoFive$ \\
$V = 3500~\kms$ &
  $\PTJWSTFullSThreeFive$ $[\PTJWSTFullSLoThreeFive, \PTJWSTFullSHiThreeFive]$ &
  $\PTSDSSFullSThreeFive$ $[\PTSDSSFullSLoThreeFive, \PTSDSSFullSHiThreeFive]$ &
  $\PTSDSSFullNThreeFive$ \\
$V = 4500~\kms$ &
  $\PTJWSTFullSFourFive$ $[\PTJWSTFullSLoFourFive, \PTJWSTFullSHiFourFive]$ &
  $\PTSDSSFullSFourFive$ $[\PTSDSSFullSLoFourFive, \PTSDSSFullSHiFourFive]$ &
  $\PTSDSSFullNFourFive$ \\
Paired $\Delta s_{2500\to 4500}$ &
  $\PTJWSTFullDeltaSigned$ $[\PTJWSTFullDeltaLo, \PTJWSTFullDeltaHi]$ &
  $\PTSDSSFullDeltaSigned$ $[\PTSDSSFullDeltaLo, \PTSDSSFullDeltaHi]$ &
  $\PTSDSSFullNFourFive$ \\
\enddata
\tablecomments{Entries give the sample standard deviation of $\log_{10}\sigV$ residuals at fixed $\log_{10}\LHa$, uncorrected for measurement error, with 95\% source-bootstrap confidence intervals in brackets. Paired increments resample sources valid at both $V = 2500$ and $4500~\kms$, so the SDSS paired $\Delta s$ differs slightly from subtracting single-window scatters with unequal $N$. Matching 95\% conditional prediction intervals for $\Delta s_{2500\to 4500}$ are $[\PTMockPresetJWSTGaussDeltaLo, \PTMockPresetJWSTGaussDeltaHi]$~dex (Gaussian, median $\PTMockPresetJWSTGaussDeltaMed$) and $[\PTMockPresetJWSTWingDeltaLo, \PTMockPresetJWSTWingDeltaHi]$~dex (core + wing, median $\PTMockPresetJWSTWingDeltaMed$) for the $\PTJWSTPresetN$ default-sideband JWST sources, and $[\PTMockFullJWSTGaussDeltaLo, \PTMockFullJWSTGaussDeltaHi]$ and $[\PTMockFullJWSTWingDeltaLo, \PTMockFullJWSTWingDeltaHi]$~dex for the $\PTJWSTFullN$ JWST sources when near-line sidebands are included. In the $\PTJWSTFullN$-source comparison, per-window cross-survey differences $s_{\mathrm{JWST}}(V) - s_{\mathrm{SDSS}}(V)$ are $\PTDiffTwoFive$~$[\PTDiffLoTwoFive, \PTDiffHiTwoFive]$, $\PTDiffThreeFive$~$[\PTDiffLoThreeFive, \PTDiffHiThreeFive]$, and $\PTDiffFourFive$~$[\PTDiffLoFourFive, \PTDiffHiFourFive]$~dex at $V = 2500$, $3500$, and $4500~\kms$.}
\end{deluxetable*}

In the $\PTJWSTPresetN$ primary JWST LRDs, widening the \Ha\ velocity half-window from $2500$ to $4500~\kms$ more than doubles the observed scatter in finite-window line dispersion at fixed \LHa, increasing $s(V)$ from $\PTJWSTPresetSTwoFive$~dex to $\PTJWSTPresetSFourFive$~dex as the median $\sigV$ grows from $\PTJWSTPresetMedSigTwoFive$ to $\PTJWSTPresetMedSigFourFive~\kms$ (Figure~\ref{fig:moments}a; Table~\ref{tab:moments}). Resampling the $\PTJWSTPresetN$ objects as pairs yields $\Delta s_{2500\to 4500} = \PTJWSTPresetDeltaSigned$~dex with a 95\% source-bootstrap confidence interval of $[\PTJWSTPresetDeltaLo, \PTJWSTPresetDeltaHi]$~dex (Figure~\ref{fig:moments}b). A comparable increase appears when moments are integrated directly on the total continuum-subtracted \Ha\ profile without narrow-line removal, giving $\Delta s_{\mathrm{tot}} = +\PTJWSTPresetRawTotDelta$~dex across the same $\PTJWSTPresetN$ sources, and within the $\PTJWSTUnflaggedN$ high-S/N LRDs without flagged Balmer absorption, which show $+\PTJWSTUnflaggedDelta$~dex with a 95\% interval of $[\PTJWSTUnflaggedDeltaLo, \PTJWSTUnflaggedDeltaHi]$~dex.

Under the same default continuum sidebands and retaining each survey's native LSF, the SDSS Type-1 AGN reference rises more gradually across the identical velocity windows, giving a paired increment of $\Delta s_{2500\to 4500} = \PTSDSSPresetDeltaSigned$~dex with a 95\% confidence interval of $[\PTSDSSPresetDeltaLo, \PTSDSSPresetDeltaHi]$~dex across the $\PTSDSSPresetNFourFive$ spectra valid in both windows. Allowing near-line continuum sidebands for detector-edge spectra (lower block of Table~\ref{tab:moments}; Appendix~\ref{app:edge}) yields consistent paired increments of $\PTJWSTFullDeltaSigned$~dex ($[\PTJWSTFullDeltaLo, \PTJWSTFullDeltaHi]$~dex) for the $\PTJWSTFullN$ JWST sources and $\PTSDSSFullDeltaSigned$~dex ($[\PTSDSSFullDeltaLo, \PTSDSSFullDeltaHi]$~dex) for $\PTSDSSFullNFourFive$ paired SDSS sources. In that $\PTJWSTFullN$-source comparison, the JWST scatter remains smaller than the SDSS scatter at $V = 2500$ and $3500~\kms$, whereas at $4500~\kms$ the cross-survey difference interval spans zero at $\PTDiffFourFive$~dex ($[\PTDiffLoFourFive, \PTDiffHiFourFive]$~dex; Table~\ref{tab:moments}), indicating no detected difference at the widest validated window without establishing statistical equivalence or a physical crossover velocity.

To test which line-profile families are consistent with the $\PTJWSTPresetN$-source increase, the lower rows of Figure~\ref{fig:moments}(b) compare $\Delta s_{2500\to 4500}$ against synthetic profiles injected into the empirical noise of the exact same $\PTJWSTPresetN$ JWST spectra (Appendix~\ref{app:injections}). Single-Gaussian inputs anchored to each object's measured \FWHMsg\ have a noiseless finite-window median increment of $\PTMockPresetJWSTGaussDeltaTrue$~dex, and propagation through empirical noise and continuum fitting yields a 95\% conditional prediction interval of $[\PTMockPresetJWSTGaussDeltaLo, \PTMockPresetJWSTGaussDeltaHi]$~dex with median $\PTMockPresetJWSTGaussDeltaMed$~dex. Core-plus-wing inputs placing $5$--$35\%$ of the broad flux into a $3\times$ wider component have a noiseless median increment of $\PTMockPresetJWSTWingDeltaTrue$~dex and a 95\% conditional prediction interval of $[\PTMockPresetJWSTWingDeltaLo, \PTMockPresetJWSTWingDeltaHi]$~dex with median $\PTMockPresetJWSTWingDeltaMed$~dex. Both prediction intervals encompass the observed $\PTJWSTPresetDeltaSigned$~dex increase, as do the corresponding $\PTJWSTFullN$-source intervals in Table~\ref{tab:moments}, so current signal-to-noise ratios do not yet separate intrinsic wing diversity from finite-window truncation and wing noise. For SDSS, both single-Gaussian-anchored mock families predict larger increments of $\PTMockPresetSDSSGaussDeltaLo$--$\PTMockPresetSDSSGaussDeltaHi$ and $\PTMockPresetSDSSWingDeltaLo$--$\PTMockPresetSDSSWingDeltaHi$~dex than the observed $\PTSDSSPresetDeltaSigned$~dex paired rise, so the SDSS spectra serve strictly as an empirical benchmark.

\subsection{Velocity-Dependent Balmer Weights}
\label{sec:res_balmer}

\begin{figure*}[t]
\centering
\includegraphics[width=\textwidth]{figures/fig2_balmer_velocity_windows.pdf}
\caption{\Hb\ becomes relatively stronger toward the wings in RUBIES-EGS-55604. Panels show velocity-binned flux ratios $R = F(\Hb)/F(\Ha)$ across five symmetric Doppler-velocity intervals relative to local systemic \OIII\ in the three highest-weight same-grating JWST LRDs: \textbf{(a)} RUBIES-EGS-55604, \textbf{(b)} RUBIES-EGS-42046, and \textbf{(c)} JADES-GN-68797. Blue circles show ratios after narrow-line subtraction and LSF matching (decomposition~A); terracotta squares show ratios after subtracting only a local continuum (decomposition~B). Both decompositions operate on the same combined grating spectrum and retain Balmer absorption. Error bars indicate $1\sigma$ pixel uncertainties, and open symbols mark intervals with $\Hb < 2\sigma$, plotted as signed ratios without positivity clipping; full numerical values and decrements $D = 1/R$ appear in Appendix Table~\ref{tab:balmer_windows}.}
\label{fig:balmer}
\end{figure*}

Figure~\ref{fig:balmer} examines how \Hb\ and \Ha\ weight the same Doppler-velocity intervals in the three highest-weight dual-line LRDs, extending the center-to-wing Balmer decrement decline found by \citet[][Section~V.3.1 and Figure~9]{Chen2026} by comparing decompositions~A and~B on identical combined grating spectra (Appendix Table~\ref{tab:balmer_windows}). Because the two decompositions differ simultaneously in continuum placement, narrow-line removal, and LSF matching, differences between them measure the combined sensitivity to spectral modeling.

In RUBIES-EGS-55604 (Figure~\ref{fig:balmer}a), which has the highest \Hb\ signal-to-noise ratio of the three targets, the continuum-only decrement $F(\Ha)/F(\Hb)$ declines from $\POneBSrcFiveFiveZeroThreeD$ at $0$--$300~\kms$ to $\POneBSrcFiveFiveOneFiveTwoFiveD$ at $1500$--$2500~\kms$ and $\POneBSrcFiveFiveTwoFiveThreeFiveD$ at $2500$--$3500~\kms$, corresponding to $R = \POneBSrcFiveFiveZeroThreeR \pm \POneBSrcFiveFiveZeroThreeRErr$ at line center and $\POneBSrcFiveFiveTwoFiveThreeFiveR \pm \POneBSrcFiveFiveTwoFiveThreeFiveRErr$ in the outer wing under decomposition~B. The same outward decline persists after narrow-line subtraction and LSF matching, with $D = \POneASrcFiveFiveZeroThreeD \to \POneASrcFiveFiveOneFiveTwoFiveD \to \POneASrcFiveFiveTwoFiveThreeFiveD$ under decomposition~A, showing that outer velocity intervals carry greater relative flux weight in \Hb\ than in \Ha\ under either spectral decomposition. By contrast, in RUBIES-EGS-42046 (Figure~\ref{fig:balmer}b) the central decrement is sensitive to narrow-line modeling, shifting from $D = \POneASrcFourTwoZeroThreeD$ in decomposition~A to $\POneBSrcFourTwoZeroThreeD$ in decomposition~B. Across $300$--$1000~\kms$, decomposition~A also shows a blue-to-red asymmetry of $D = \POneASrcFourTwoThreeTenBlueD$ versus $\POneASrcFourTwoThreeTenRedD$ overlapping the blueshifted Balmer absorber reported in Table~2 of \citet{Chen2026}, while \Hb\ falls below $2\sigma$ at $2500$--$3500~\kms$. In JADES-GN-68797 (Figure~\ref{fig:balmer}c), continuum-only subtraction leaves \Hb\ below $2\sigma$ in every bin beyond $1000~\kms$.

Comparing these velocity-binned profiles with the $\POneTmplAllThirtyThreeN$-source inverse-variance-weighted template ratio of $\POneTmplAllThirtyThreeRatio$ (95\% conditional interval $[\POneTmplAllThirtyThreeLo, \POneTmplAllThirtyThreeHi]$; Appendix~\ref{app:template_weights}) clarifies why the two statistics differ. The template fit normalizes each object's broad-\Ha\ profile to \Hb\ across $300 \le |v| \le 1000~\kms$ relative to the broad-\Ha\ centroid and measures an inverse-variance-weighted wing amplitude across $1500 \le |v| \le 4500~\kms$, which is dominated by the brighter $1500$--$2500~\kms$ shoulder, whereas Figure~\ref{fig:balmer} evaluates each velocity bin directly down to $0$--$300~\kms$ relative to systemic \OIII. A near-unity wing-to-shoulder template ratio across the $\POneTmplTopThreeN$ dominant objects, which supply $\POneTmplTopThreeWtPct\%$ of the total template weight, therefore coexists naturally with the center-to-wing gradient in RUBIES-EGS-55604, while the remaining $\POneTmplOtherThirtyN$ fainter targets carry $\POneTmplOtherThirtyWtPct\%$ of the weight and remain poorly constrained at $\POneTmplOtherThirtyRatio$ with a 95\% interval of $[\POneTmplOtherThirtyLo, \POneTmplOtherThirtyHi]$ spanning both zero and unity.

\subsection{Integrated-Width Population Context}
\label{sec:res_context}

\begin{figure*}[t]
\centering
\includegraphics[width=\textwidth]{figures/fig3_three_survey_width_luminosity.pdf}
\caption{Integrated broad-\Ha\ single-Gaussian widths (\FWHMsg) in three independently selected samples on common axes, with uniform OLS residual scatter $s$ at fixed observed \LHa\ listed above each panel. Gold shading marks the shared luminosity interval $\log_{10}(\LHa/\mathrm{erg\,s^{-1}}) = 43.0$--$43.7$; outer black rings in panel~\textbf{(a)} identify its $\JWSTBridgeN$ JWST sources. \textbf{(a)} JWST/NIRSpec LRDs ($N = \JWSTCleanN$ continuum-clean baseline, distinguishing the $\PTJWSTPresetN$ primary default-sideband window sources in filled crimson from the remaining $\JWSTOtherCleanN$ continuum-clean sources). \textbf{(b)} Quality-controlled SDSS DR7 Type-1 AGNs ($N = \SDSSCleanN$). \textbf{(c)} Euclid DR1 primary broad-\Ha\ sample ($N = \EuclidMainN$, $19.0 < H_{\mathrm{aper}} < \EuclidHLimit$).}
\label{fig:three_survey}
\end{figure*}

Figure~\ref{fig:three_survey} places integrated \FWHMsg\ at fixed observed \LHa\ on common axes across all three surveys using the identical OLS residual standard deviation. The $\JWSTCleanN$ continuum-clean JWST LRDs span a much narrower range of integrated widths, with $s = \JWSTCleanSGScatterUniform$~dex across all $\JWSTCleanN$ objects and $\PTJWSTPresetSGScatter$~dex across the $\PTJWSTPresetN$ primary window targets, than either the $\SDSSCleanN$ SDSS Type-1 AGNs at $s = \SDSSCleanSGScatterUniform$~dex or the $\EuclidMainN$ Euclid DR1 AGNs at $s = \EuclidMainSGScatterUniform$~dex. Independently, hierarchical modeling of the $\EuclidMainN$ Euclid DR1 AGNs conditioned on total rest-optical continuum luminosity and corrected for flux-limit selection yields an intrinsic width scatter of $\sigma_w = \EuclidMainSigmaW$~dex with a 95\% credible interval of $[\EuclidMainSigmaWLowerNinetyFive, \EuclidMainSigmaWUpperNinetyFive]$~dex. Because only $\JWSTBridgeN$ JWST sources fall inside the shared $\log_{10}(\LHa/\mathrm{erg\,s^{-1}}) = 43.0$--$43.7$ band alongside $\SDSSBridgeN$ SDSS and $\EuclidBridgeN$ Euclid AGNs, and a $\BridgeJWSTEuclidContinuumGap$~dex continuum luminosity gap separates the brightest JWST and faintest Euclid overlap objects, the three surveys provide a population reference for integrated \Ha\ widths rather than a test of redshift evolution at fixed luminosity.

\section{Discussion}
\label{sec:discussion}

Our two profile analyses probe complementary observables: population-level scatter in the \Ha\ finite-window line dispersion (Figure~\ref{fig:moments}) and single-object velocity-binned \Hb/\Ha\ ratios (Figure~\ref{fig:balmer}). Together, they clarify what an integrated LRD line width measures. Within a narrow symmetric window where flux density varies modestly across $[-V, +V]$, the finite-window line dispersion approaches $\sigV = V/\sqrt{3}$ (Appendix~\ref{app:exact_int}), compressing differences among profiles of distinct overall breadth. As $V$ widens from $2500$ to $4500~\kms$, outer pixels gain $(v - \bar{v}_V)^2$ weight and the observed scatter across the same $\PTJWSTPresetN$ LRDs more than doubles from $\PTJWSTPresetSTwoFive$ to $\PTJWSTPresetSFourFive$~dex. Because Gaussian and core-plus-wing injections into the empirical noise of those $\PTJWSTPresetN$ spectra both accommodate the observed $\PTJWSTPresetDeltaSigned$~dex rise, the present data establish how the measured scatter depends on velocity window without isolating a noise-free ``line-core'' invariant or separating intrinsic wing diversity from truncation and noise.

On the individual-source level, a uniform foreground dust screen attenuates $F(\Hb)/F(\Ha)$ by nearly the same factor across all Doppler velocities of a single line; while dust can shift the integrated Balmer decrement, producing a center-to-wing gradient requires velocity-dependent intrinsic emission or differential radiative transfer. In RUBIES-EGS-55604, outer Doppler intervals carry a higher \Hb/\Ha\ flux ratio under either spectral decomposition, corroborating \citet{Chen2026}. Physically, higher-$|v|$ emission can arise in stratified broad-line gas where density and $n=2$ optical depth vary across emission zones \citep{OsterbrockFerland2006,InayoshiMaiolino2025,Madau2026b,Scholtz2026}, alongside velocity-selective Balmer absorption in dense gas \citep{deGraaffRuby2025,DEugenio2026,Matthee2026} or frequency redistribution and differential escape through optically thick envelopes \citep[][and the review by \citealt{InayoshiHo}]{Laor2006,Brazzini2025,Rusakov2026,Chang2025,NaiduBHStar2026,deGraaff2026}. Because scattering and absorption can decouple observed Doppler shift $v$ from bulk orbital speed, velocity-binned ratios alone neither map $v$ uniquely onto radius nor establish whether \Ha\ and \Hb\ yield discrepant black hole masses.

Viewed alongside the broader integrated-width distributions of SDSS and Euclid DR1 AGNs in Figure~\ref{fig:three_survey}, these findings set a direct forward-modeling requirement. Without independent black hole masses or reverberation lags, our spectra do not constrain the numerical value of a universal virial factor \fvir. Instead, any dynamical, optical-depth, or scattering model that uses a finite-window line dispersion or integrated width to infer LRD black hole masses or kinematics must evaluate its predictions over the \emph{same Balmer transition, width definition, and velocity interval} as the observations, reproducing both the window-dependent scatter $s(V)$ in Figure~\ref{fig:moments} and the velocity-binned Balmer ratios in bright dual-line targets such as RUBIES-EGS-55604. Deeper NIRSpec grating spectra will be needed to validate \Ha\ moments beyond $4500~\kms$, extend resolved \Hb/\Ha\ profiles to fainter LRDs, and test width-scaled windows once boundary uncertainties are propagated.

\section{Conclusions}
\label{sec:conclusions}

Holding the objects, width definition, and broad-\Ha\ luminosities fixed across velocity intervals yields two main empirical conclusions:
\begin{enumerate}
\item Across $\PTJWSTPresetN$ high-S/N JWST LRDs with default continuum sidebands, the observed scatter in broad-\Ha\ finite-window line dispersion \sigV\ at fixed \LHa\ rises from $\PTJWSTPresetSTwoFive$~dex at $V = 2500~\kms$ to $\PTJWSTPresetSFourFive$~dex at $V = 4500~\kms$, a paired increase of $\PTJWSTPresetDeltaSigned$~dex with 95\% bootstrap interval $[\PTJWSTPresetDeltaLo, \PTJWSTPresetDeltaHi]$~dex (Table~\ref{tab:moments}), versus $\PTSDSSPresetDeltaSigned$~dex ($[\PTSDSSPresetDeltaLo, \PTSDSSPresetDeltaHi]$~dex) in $\PTSDSSPresetNFourFive$ SDSS Type-1 AGNs and integrated-width scatters of $\JWSTCleanSGScatterUniform$, $\SDSSCleanSGScatterUniform$, and $\EuclidMainSGScatterUniform$~dex across JWST ($N = \JWSTCleanN$), SDSS ($N = \SDSSCleanN$), and Euclid DR1 ($N = \EuclidMainN$). Synthetic profile injections into the empirical noise of the same $\PTJWSTPresetN$ JWST spectra show that single-Gaussian and core-plus-wing models both span the observed $\Delta s_{2500\to 4500}$, leaving intrinsic wing diversity, truncation, and wing noise unseparated.
\item In RUBIES-EGS-55604, outer velocity intervals carry greater relative \Hb\ flux than \Ha\ under both narrow-line-subtracted and continuum-only decompositions of the same combined grating spectrum, with $F(\Ha)/F(\Hb)$ declining from $\POneBSrcFiveFiveZeroThreeD$ to $\POneBSrcFiveFiveOneFiveTwoFiveD$ and $\POneBSrcFiveFiveTwoFiveThreeFiveD$ from center to wing after continuum subtraction alone. Comparisons between LRD broad-line models and observations must therefore match both the Balmer transition and the velocity interval included in the finite-window line dispersion.
\end{enumerate}

\section*{Acknowledgments}
The author is grateful to his parents for their constant support. This research received no external grant funding.

Some of the data products presented herein were retrieved from the Dawn JWST Archive (DJA). DJA is an initiative of the Cosmic Dawn Center (DAWN), which is funded by the Danish National Research Foundation under grant DNRF140. This work is based on observations made with the NASA/ESA/CSA James Webb Space Telescope. The reduced spectra used here were retrieved from the DAWN JWST Archive; the underlying JWST observations are archived in the Mikulski Archive for Space Telescopes at the Space Telescope Science Institute, which is operated by the Association of Universities for Research in Astronomy, Inc., under NASA contract NAS5-03127 for JWST. These observations are associated with programs \#1181 (PI D.~J. Eisenstein), \#1210 (PI N. L\"utzgendorf), \#1286 (PI N. L\"utzgendorf), \#1345 (PI S.~L. Finkelstein), \#4106 (PI E. Nelson), \#4233 (PI A. de Graaff), \#5718 (PI D.~D. Kocevski), and \#8047 (PI B. Wang). The author acknowledges the teams of programs \#1345 and \#4233 for developing their observing programs with zero-exclusive-access periods.

Funding for the SDSS and SDSS-II has been provided by the Alfred P. Sloan Foundation, the Participating Institutions, the National Science Foundation, the U.S. Department of Energy, the National Aeronautics and Space Administration, the Japanese Monbukagakusho, the Max Planck Society, and the Higher Education Funding Council for England. The SDSS Web Site is \url{http://www.sdss.org/}. The SDSS is managed by the Astrophysical Research Consortium for the Participating Institutions.

The Participating Institutions are the American Museum of Natural History, Astrophysical Institute Potsdam, University of Basel, University of Cambridge, Case Western Reserve University, University of Chicago, Drexel University, Fermilab, the Institute for Advanced Study, the Japan Participation Group, Johns Hopkins University, the Joint Institute for Nuclear Astrophysics, the Kavli Institute for Particle Astrophysics and Cosmology, the Korean Scientist Group, the Chinese Academy of Sciences (LAMOST), Los Alamos National Laboratory, the Max-Planck-Institute for Astronomy (MPIA), the Max-Planck-Institute for Astrophysics (MPA), New Mexico State University, Ohio State University, University of Pittsburgh, University of Portsmouth, Princeton University, the United States Naval Observatory, and the University of Washington.

This work has made use of the Euclid DR1 data from the \emph{Euclid} mission of the European Space Agency (ESA), 2026, \url{https://doi.org/10.57780/esa-d079c51}, processed by the \emph{Euclid} Consortium (\url{https://www.euclid-ec.org/}).

Reproduction scripts, derived catalogs, and figure-generation code for this study are archived at \url{https://github.com/NianSi-AI/lrd-balmer-velocity-windows} and on Zenodo at \dataset[doi:10.5281/zenodo.24426322]{https://doi.org/10.5281/zenodo.24426322}.

AI disclosure: Gemini (Google), Claude (Anthropic), Codex (OpenAI), and DeepSeek were used under the author's direction to audit physical and statistical arguments, write and execute analysis and Matplotlib figure-generation code, cross-check literature citations, and refine prose from the author's drafts. All text, numbers, figures, and references were verified by the author, who bears full responsibility for the content.

\section*{Author Contributions}
H.W. conceived and designed the study, performed the data analysis and modeling, and wrote the manuscript.

\facilities{JWST (NIRSpec), Sloan, Euclid (NISP, VIS)}

\software{\texttt{msaexp} \citep{msaexp}, NumPy \citep{harris2020array}, SciPy \citep{2020SciPy-NMeth}, Matplotlib \citep{Hunter:2007}, Astropy \citep{astropy:2013,astropy:2018,astropy:2022}}

\clearpage
\appendix
\restartappendixnumbering
\makeatletter
\@addtoreset{figure}{section}
\@addtoreset{table}{section}
\makeatother
\renewcommand{\theHfigure}{\Alph{section}\arabic{figure}}
\renewcommand{\theHtable}{\Alph{section}\arabic{table}}
\twocolumngrid

\section{Finite-Window Integration, Recovery Calibration, and Near-Line Sidebands}
\label{app:calibration}

\begin{figure*}[t]
\centering
\includegraphics[width=0.92\textwidth]{figures/figA1_moment_calibration.pdf}
\caption{Recovery calibration and extended-window diagnostics for the \Ha\ finite-window line dispersion. \textbf{(a)} 68th-percentile $|\Delta\log_{10}\sigV|$ across $V = 1500$--$6000~\kms$. Open symbols fail at least one per-profile qualification threshold ($\ge 85\%$ valid recoveries, $|\text{median log bias}| \le 0.03$~dex, or 68th-percentile $|\Delta\log_{10}\sigV| \le 0.08$~dex; the dashed horizontal line marks only the $0.08$~dex error limit, so JWST core-plus-wing at $V = 6000~\kms$ is open owing to its $\PTCalJWSTWingSixZeroUsablePct\%$ recovery rate). \textbf{(b)} Extended-window observed scatter diagnostic showing validated $V \le 4500~\kms$ curves alongside disqualified JWST windows at $V = 5500$ and $6000~\kms$ (open circles).}
\label{fig:app_cal}
\end{figure*}

\subsection{Positive-Overlap Pixel Quadrature and Flat-Window Limit}
\label{app:exact_int}

On each native wavelength grid $\{\lambda_i\}$, midpoint edges $\lambda_{i,\pm} = \tfrac{1}{2}(\lambda_i + \lambda_{i\pm 1})$, extrapolated symmetrically at spectrum endpoints, convert to rest-frame velocity edges $[v_{i,-}, v_{i,+}]$. In the \Ha\ moment scan, each pixel edge is further restricted to lie within $0.55$ times the median pixel spacing of $v_i$ so that masked or missing wavelength intervals are never bridged. Intersecting each valid pixel with $[v_{\mathrm{lo}}, v_{\mathrm{hi}}]$ and retaining only positive overlaps via $a_i = \max(v_{i,-}, v_{\mathrm{lo}})$ and $b_i = \max(a_i, \min(v_{i,+}, v_{\mathrm{hi}}))$ gives exact piecewise-constant integrals for window flux $F(V) = \sum_i f_i (b_i - a_i)$, centroid $\bar{v}_V = [2 F(V)]^{-1}\sum_i f_i (b_i^2 - a_i^2)$, and finite-window line dispersion
\begin{equation}
\sigV = \left[ \frac{1}{3 F(V)} \sum_i f_i \left( (b_i - \bar{v}_V)^3 - (a_i - \bar{v}_V)^3 \right) \right]^{1/2}.
\label{eq:app_sigma_v}
\end{equation}
For constant flux $f_i = f_0$ across $[-V, +V]$, Equation~(\ref{eq:app_sigma_v}) yields $F(V) = 2 f_0 V$, $\bar{v}_V = 0$, and $\sigV = V/\sqrt{3}$. Because finite $V$ truncates broader lines more strongly than narrower lines, the noiseless $\PTJWSTPresetN$-source Gaussian mock scatter at $V = 2500~\kms$ ($\PTMockPresetJWSTGaussTrueTwoFive$~dex for JWST and $\PTMockPresetSDSSGaussTrueTwoFive$~dex for SDSS) lies below its untruncated full-profile counterpart of $\PTMockPresetJWSTGaussTrueFull$ and $\PTMockPresetSDSSGaussTrueFull$~dex, respectively.

\subsection{Window Recovery Calibration and Empirical-Noise Mock Ensembles}
\label{app:injections}

\begin{deluxetable*}{lcccccc}
\tabletypesize{\scriptsize}
\tabcolsep=4.5pt
\tablewidth{0pt}
\tablecaption{Injection-Recovery Window Calibration Across Tested Velocity Windows\label{tab:calibration}}
\tablehead{
\colhead{Window $V$} &
\multicolumn{3}{c}{Gaussian Input} &
\multicolumn{3}{c}{Core + Wing Input ($5$--$35\%$ Wing Flux)} \\
\cmidrule(lr){2-4} \cmidrule(lr){5-7}
\colhead{(\kms)} &
\colhead{Recovered (\%)} &
\colhead{Med.\ Bias (dex)} &
\colhead{68th Err.\ (dex)} &
\colhead{Recovered (\%)} &
\colhead{Med.\ Bias (dex)} &
\colhead{68th Err.\ (dex)}
}
\startdata
\cutinhead{JWST, $\PTCalJWSTAttemptedTargetsN$ targets attempted (default sidebands)}
1500 & $\PTCalJWSTGaussOneFiveUsablePct$ & $\PTCalJWSTGaussOneFiveBias$ & $\PTCalJWSTGaussOneFiveErrSixEight$ & $\PTCalJWSTWingOneFiveUsablePct$ & $\PTCalJWSTWingOneFiveBias$ & $\PTCalJWSTWingOneFiveErrSixEight$ \\
2500 & $\PTCalJWSTGaussTwoFiveUsablePct$ & $\PTCalJWSTGaussTwoFiveBias$ & $\PTCalJWSTGaussTwoFiveErrSixEight$ & $\PTCalJWSTWingTwoFiveUsablePct$ & $\PTCalJWSTWingTwoFiveBias$ & $\PTCalJWSTWingTwoFiveErrSixEight$ \\
3500 & $\PTCalJWSTGaussThreeFiveUsablePct$ & $\PTCalJWSTGaussThreeFiveBias$ & $\PTCalJWSTGaussThreeFiveErrSixEight$ & $\PTCalJWSTWingThreeFiveUsablePct$ & $\PTCalJWSTWingThreeFiveBias$ & $\PTCalJWSTWingThreeFiveErrSixEight$ \\
4500 & $\PTCalJWSTGaussFourFiveUsablePct$ & $\PTCalJWSTGaussFourFiveBias$ & $\PTCalJWSTGaussFourFiveErrSixEight$ & $\PTCalJWSTWingFourFiveUsablePct$ & $\PTCalJWSTWingFourFiveBias$ & $\PTCalJWSTWingFourFiveErrSixEight$ \\
5500\tablenotemark{a} & $\PTCalJWSTGaussFiveFiveUsablePct^{\ast}$ & $\PTCalJWSTGaussFiveFiveBias$ & $\PTCalJWSTGaussFiveFiveErrSixEight^{\ast}$ & $\PTCalJWSTWingFiveFiveUsablePct$ & $\PTCalJWSTWingFiveFiveBias$ & $\PTCalJWSTWingFiveFiveErrSixEight$ \\
6000\tablenotemark{a} & $\PTCalJWSTGaussSixZeroUsablePct^{\ast}$ & $\PTCalJWSTGaussSixZeroBias$ & $\PTCalJWSTGaussSixZeroErrSixEight^{\ast}$ & $\PTCalJWSTWingSixZeroUsablePct^{\ast}$ & $\PTCalJWSTWingSixZeroBias$ & $\PTCalJWSTWingSixZeroErrSixEight$ \\
\cutinhead{SDSS, $\PTCalSDSSAttemptedTargetsN$ targets attempted (default sidebands)}
1500 & $\PTCalSDSSGaussOneFiveUsablePct$ & $\PTCalSDSSGaussOneFiveBias$ & $\PTCalSDSSGaussOneFiveErrSixEight$ & $\PTCalSDSSWingOneFiveUsablePct$ & $\PTCalSDSSWingOneFiveBias$ & $\PTCalSDSSWingOneFiveErrSixEight$ \\
2500 & $\PTCalSDSSGaussTwoFiveUsablePct$ & $\PTCalSDSSGaussTwoFiveBias$ & $\PTCalSDSSGaussTwoFiveErrSixEight$ & $\PTCalSDSSWingTwoFiveUsablePct$ & $\PTCalSDSSWingTwoFiveBias$ & $\PTCalSDSSWingTwoFiveErrSixEight$ \\
3500 & $\PTCalSDSSGaussThreeFiveUsablePct$ & $\PTCalSDSSGaussThreeFiveBias$ & $\PTCalSDSSGaussThreeFiveErrSixEight$ & $\PTCalSDSSWingThreeFiveUsablePct$ & $\PTCalSDSSWingThreeFiveBias$ & $\PTCalSDSSWingThreeFiveErrSixEight$ \\
4500 & $\PTCalSDSSGaussFourFiveUsablePct$ & $\PTCalSDSSGaussFourFiveBias$ & $\PTCalSDSSGaussFourFiveErrSixEight$ & $\PTCalSDSSWingFourFiveUsablePct$ & $\PTCalSDSSWingFourFiveBias$ & $\PTCalSDSSWingFourFiveErrSixEight$ \\
5500 & $\PTCalSDSSGaussFiveFiveUsablePct$ & $\PTCalSDSSGaussFiveFiveBias$ & $\PTCalSDSSGaussFiveFiveErrSixEight$ & $\PTCalSDSSWingFiveFiveUsablePct$ & $\PTCalSDSSWingFiveFiveBias$ & $\PTCalSDSSWingFiveFiveErrSixEight$ \\
6000 & $\PTCalSDSSGaussSixZeroUsablePct$ & $\PTCalSDSSGaussSixZeroBias$ & $\PTCalSDSSGaussSixZeroErrSixEight$ & $\PTCalSDSSWingSixZeroUsablePct$ & $\PTCalSDSSWingSixZeroBias$ & $\PTCalSDSSWingSixZeroErrSixEight$ \\
\enddata
\tablecomments{Calibration tests injected $10$ empirical four-pixel block-noise realizations per attempted target and profile class under default sidebands, so recovered percentages divide valid recoveries by all attempted realizations in each cell (for example, $\PTCalJWSTGaussOneFiveOkTrials/\PTCalJWSTGaussOneFiveTotTrials = \PTCalJWSTGaussOneFiveUsablePct\%$ for JWST and $\PTCalSDSSGaussOneFiveOkTrials/\PTCalSDSSGaussOneFiveTotTrials = \PTCalSDSSGaussOneFiveUsablePct\%$ for SDSS at $V = 1500~\kms$ Gaussian), including spectra that lack default sidebands. Joint window qualification requires both profile classes to achieve $\ge 85\%$ valid recoveries, $|\text{median log bias}| \le 0.03$~dex, and 68th-percentile $|\Delta\log_{10}\sigV| \le 0.08$~dex; asterisks ($^{\ast}$) mark individual threshold violations.}
\tablenotetext{a}{Loses joint JWST window qualification: $V = 5500~\kms$ fails the Gaussian recovery rate and error thresholds, and $V = 6000~\kms$ fails the Gaussian recovery rate and error thresholds as well as the core-plus-wing recovery rate. Across all $\JWSTCleanN$ continuum-clean JWST targets ($\PTJWSTFortyNFourFive$ valid at $4500~\kms$, with observed $s = \PTJWSTFortySTwoFive \to \PTJWSTFortySFourFive$~dex), Gaussian recovery already falls short at $V = 4500~\kms$ ($\PTCalJWSTFortyGaussFourFiveUsablePct\%$ recovered, $\PTCalJWSTFortyGaussFourFiveErrSixEight$~dex 68th-percentile error).}
\end{deluxetable*}

We established the window qualification thresholds prior to measuring the observed finite-window moments (Figure~\ref{fig:app_cal}; Table~\ref{tab:calibration}) by injecting $10$ four-pixel block-noise realizations per attempted target across $\PTCalJWSTAttemptedTargetsN$ high-S/N JWST and $\PTCalSDSSAttemptedTargetsN$ SDSS spectra, including synthetic broad \Ha, narrow \Ha+\NII, linear continua, and blueshifted absorption in a subset of realizations. To construct the conditional prediction intervals for $\Delta s_{2500\to 4500}$, we then generated $60$ empirical-noise realizations per JWST source and $12$ per SDSS source for each profile class, fixing each object's intrinsic core FWHM to its measured \FWHMsg\ and including blueshifted absorption in a subset of realizations prior to LSF convolution. Drawing one realization recovered in both windows per source preserves paired pixel noise across $4000$ synthetic sample draws, yielding the $\PTJWSTPresetN$-source default-sideband prediction intervals in Figure~\ref{fig:moments}(b) and the $\PTJWSTFullN$-source intervals in Table~\ref{tab:moments}.

\subsection{Near-Line Continuum Sidebands for Detector-Edge Spectra}
\label{app:edge}

Following the primary default-sideband analysis, we tested closer continuum windows at $[+5600, +6300]~\kms$ on the red side and $[-6300, -5400]~\kms$ on the blue side for spectra truncated near detector edges: RUBIES-EGS-50052 and JADES-GS-38562 in JWST, together with the additional SDSS spectra entering the $\PTSDSSFullNFourFive$--$\PTSDSSFullNOneFive$ set in Table~\ref{tab:moments}. At $V = 4500~\kms$, RUBIES-EGS-50052 achieves $\PTEdgeFiveZeroZeroFiveTwoGaussUsablePct\%$ valid recoveries with a $\PTEdgeFiveZeroZeroFiveTwoGaussErrSixEight$~dex 68th-percentile error for Gaussian inputs and $\PTEdgeFiveZeroZeroFiveTwoWingUsablePct\%$ with a $\PTEdgeFiveZeroZeroFiveTwoWingErrSixEight$~dex error for core-plus-wing inputs. By contrast, JADES-GS-38562 is noisier near the detector cutoff, reaching $\PTEdgeThreeEightFiveSixTwoGaussUsablePct\%$ valid recoveries and $\PTEdgeThreeEightFiveSixTwoGaussErrSixEight$~dex error for Gaussian inputs and $\PTEdgeThreeEightFiveSixTwoWingUsablePct\%$ and $\PTEdgeThreeEightFiveSixTwoWingErrSixEight$~dex for core-plus-wing inputs, motivating the $\PTJWSTPresetN$ default-sideband objects as our primary sample.

\section{Velocity-Binned Balmer Ratios and Template Weight Breakdown}
\label{app:template_weights}

Table~\ref{tab:balmer_windows} lists the five-interval $R = F(\Hb)/F(\Ha)$ and $D = 1/R$ measurements shown in Figure~\ref{fig:balmer}. Figure~\ref{fig:app_template} shows how individual targets weight the $\POneTmplAllThirtyThreeN$-source same-grating template ratio discussed in Section~\ref{sec:res_balmer}, using inverse-variance weights calibrated on $4096$ empirical block-noise realizations per source and evaluated on an independent set of $4096$ realizations. Because inverse-variance weights scale quadratically with line S/N, RUBIES-EGS-55604, RUBIES-EGS-42046, and JADES-GN-68797 supply $\POneTmplTopThreeWtPct\%$ of the total weight, yielding a wing-to-shoulder ratio of $\POneTmplTopThreeRatio$ with a 95\% interval of $[\POneTmplTopThreeLo, \POneTmplTopThreeHi]$. The remaining $\POneTmplOtherThirtyN$ fainter objects carry $\POneTmplOtherThirtyWtPct\%$ of the weight and give $\POneTmplOtherThirtyRatio$ with a 95\% interval of $[\POneTmplOtherThirtyLo, \POneTmplOtherThirtyHi]$, spanning both zero and unity because noisy \Hb\ amplitudes are retained without positivity clipping.

\begin{figure*}[t]
\centering
\includegraphics[width=0.52\textwidth]{figures/figB1_template_weights.pdf}
\caption{\Hb\ template amplitude in the $1500$--$4500~\kms$ wing band relative to the $300$--$1000~\kms$ normalization band, with 95\% pixel-noise intervals conditional on fixed \Ha\ profiles, across all $\POneTmplAllThirtyThreeN$ same-grating sources ($\POneTmplAllThirtyThreeRatio$, $[\POneTmplAllThirtyThreeLo, \POneTmplAllThirtyThreeHi]$), the top $\POneTmplTopThreeN$ sources shown in Figure~\ref{fig:balmer} ($\POneTmplTopThreeWtPct\%$ weight; $\POneTmplTopThreeRatio$, $[\POneTmplTopThreeLo, \POneTmplTopThreeHi]$), and the remaining $\POneTmplOtherThirtyN$ sources ($\POneTmplOtherThirtyWtPct\%$ weight; $\POneTmplOtherThirtyRatio$, $[\POneTmplOtherThirtyLo, \POneTmplOtherThirtyHi]$).}
\label{fig:app_template}
\end{figure*}

\begin{deluxetable*}{lcccccc}
\tabletypesize{\scriptsize}
\tablewidth{0pt}
\tablecaption{Velocity-Binned \Hb/\Ha\ Flux Ratios $R$ and Equivalent Decrements $D$ in Three High-Weight Same-Grating LRDs\label{tab:balmer_windows}}
\tablehead{
\colhead{$|v|$ Interval} &
\multicolumn{2}{c}{RUBIES-EGS-55604} &
\multicolumn{2}{c}{RUBIES-EGS-42046} &
\multicolumn{2}{c}{JADES-GN-68797} \\
\cmidrule(lr){2-3} \cmidrule(lr){4-5} \cmidrule(lr){6-7}
\colhead{(\kms)} &
\colhead{Narrow sub.\ (A)} &
\colhead{Cont.\ only (B)} &
\colhead{Narrow sub.\ (A)} &
\colhead{Cont.\ only (B)} &
\colhead{Narrow sub.\ (A)} &
\colhead{Cont.\ only (B)}
}
\startdata
$0$--$300$ &
  $\POneASrcFiveFiveZeroThreeR \pm \POneASrcFiveFiveZeroThreeRErr$ ($\POneASrcFiveFiveZeroThreeD$) &
  $\POneBSrcFiveFiveZeroThreeR \pm \POneBSrcFiveFiveZeroThreeRErr$ ($\POneBSrcFiveFiveZeroThreeD$) &
  $\POneASrcFourTwoZeroThreeR \pm \POneASrcFourTwoZeroThreeRErr$ ($\POneASrcFourTwoZeroThreeD$) &
  $\POneBSrcFourTwoZeroThreeR \pm \POneBSrcFourTwoZeroThreeRErr$ ($\POneBSrcFourTwoZeroThreeD$) &
  $\POneASrcSixEightZeroThreeR \pm \POneASrcSixEightZeroThreeRErr$ ($\POneASrcSixEightZeroThreeD$) &
  $\POneBSrcSixEightZeroThreeR \pm \POneBSrcSixEightZeroThreeRErr$ ($\POneBSrcSixEightZeroThreeD$) \\
$300$--$1000$ &
  $\POneASrcFiveFiveThreeTenR \pm \POneASrcFiveFiveThreeTenRErr$ ($\POneASrcFiveFiveThreeTenD$) &
  $\POneBSrcFiveFiveThreeTenR \pm \POneBSrcFiveFiveThreeTenRErr$ ($\POneBSrcFiveFiveThreeTenD$) &
  $\POneASrcFourTwoThreeTenR \pm \POneASrcFourTwoThreeTenRErr$ ($\POneASrcFourTwoThreeTenD$) &
  $\POneBSrcFourTwoThreeTenR \pm \POneBSrcFourTwoThreeTenRErr$ ($\POneBSrcFourTwoThreeTenD$) &
  $\POneASrcSixEightThreeTenR \pm \POneASrcSixEightThreeTenRErr$ ($\POneASrcSixEightThreeTenD$) &
  $\POneBSrcSixEightThreeTenR \pm \POneBSrcSixEightThreeTenRErr$ ($\POneBSrcSixEightThreeTenD$) \\
$1000$--$1500$ &
  $\POneASrcFiveFiveTenOneFiveR \pm \POneASrcFiveFiveTenOneFiveRErr$ ($\POneASrcFiveFiveTenOneFiveD$) &
  $\POneBSrcFiveFiveTenOneFiveR \pm \POneBSrcFiveFiveTenOneFiveRErr$ ($\POneBSrcFiveFiveTenOneFiveD$) &
  $\POneASrcFourTwoTenOneFiveR \pm \POneASrcFourTwoTenOneFiveRErr$ ($\POneASrcFourTwoTenOneFiveD$) &
  $\POneBSrcFourTwoTenOneFiveR \pm \POneBSrcFourTwoTenOneFiveRErr$ ($\POneBSrcFourTwoTenOneFiveD$) &
  $\POneASrcSixEightTenOneFiveR \pm \POneASrcSixEightTenOneFiveRErr$ ($\POneASrcSixEightTenOneFiveD$) &
  $\POneBSrcSixEightTenOneFiveR \pm \POneBSrcSixEightTenOneFiveRErr$ ($<2\sigma$) \\
$1500$--$2500$ &
  $\POneASrcFiveFiveOneFiveTwoFiveR \pm \POneASrcFiveFiveOneFiveTwoFiveRErr$ ($\POneASrcFiveFiveOneFiveTwoFiveD$) &
  $\POneBSrcFiveFiveOneFiveTwoFiveR \pm \POneBSrcFiveFiveOneFiveTwoFiveRErr$ ($\POneBSrcFiveFiveOneFiveTwoFiveD$) &
  $\POneASrcFourTwoOneFiveTwoFiveR \pm \POneASrcFourTwoOneFiveTwoFiveRErr$ ($\POneASrcFourTwoOneFiveTwoFiveD$) &
  $\POneBSrcFourTwoOneFiveTwoFiveR \pm \POneBSrcFourTwoOneFiveTwoFiveRErr$ ($\POneBSrcFourTwoOneFiveTwoFiveD$) &
  $\POneASrcSixEightOneFiveTwoFiveR \pm \POneASrcSixEightOneFiveTwoFiveRErr$ ($\POneASrcSixEightOneFiveTwoFiveD$) &
  $\POneBSrcSixEightOneFiveTwoFiveR \pm \POneBSrcSixEightOneFiveTwoFiveRErr$ ($<2\sigma$) \\
$2500$--$3500$ &
  $\POneASrcFiveFiveTwoFiveThreeFiveR \pm \POneASrcFiveFiveTwoFiveThreeFiveRErr$ ($\POneASrcFiveFiveTwoFiveThreeFiveD$) &
  $\POneBSrcFiveFiveTwoFiveThreeFiveR \pm \POneBSrcFiveFiveTwoFiveThreeFiveRErr$ ($\POneBSrcFiveFiveTwoFiveThreeFiveD$) &
  $\POneASrcFourTwoTwoFiveThreeFiveR \pm \POneASrcFourTwoTwoFiveThreeFiveRErr$ ($<2\sigma$) &
  $\POneBSrcFourTwoTwoFiveThreeFiveR \pm \POneBSrcFourTwoTwoFiveThreeFiveRErr$ ($<2\sigma$) &
  $\POneASrcSixEightTwoFiveThreeFiveR \pm \POneASrcSixEightTwoFiveThreeFiveRErr$ ($<2\sigma$) &
  $\POneBSrcSixEightTwoFiveThreeFiveR \pm \POneBSrcSixEightTwoFiveThreeFiveRErr$ ($<2\sigma$) \\
\enddata
\tablecomments{Entries list $R \equiv F(\Hb)/F(\Ha) \pm 1\sigma$ integrated over symmetric native-pixel overlaps relative to local systemic \OIII, propagating pixel uncertainties plus local continuum-fit covariance in decomposition~B. Equivalent decrements $D = 1/R$ appear in parentheses only where integrated $\Hb \ge 2\sigma$.}
\end{deluxetable*}

\onecolumngrid
\clearpage
\raggedbottom
\twocolumngrid
\bibliography{references}
\bibliographystyle{aasjournal}

\end{document}
